\documentclass{article}
\usepackage[utf8]{inputenc}
\usepackage{romannum}
\usepackage{amsfonts}
\usepackage{amssymb}
\usepackage{fancyhdr}
\usepackage{graphicx}
\usepackage{t1enc}
\usepackage{pdfpages}
\usepackage[magyar]{babel}
\usepackage[utf8]{inputenc}
\usepackage{amsmath}
\usepackage{mathtools}
\usepackage{pdfpages}
\usepackage{ marvosym } 
\usepackage{wrapfig}
\usepackage{hyperref}
\usepackage{pgfplots}

\pgfplotsset{width=\textwidth,compat=newest}
\hypersetup{
    colorlinks,
    citecolor=black,
    filecolor=black,
    linkcolor=black,
    urlcolor=black
}
\usepackage{romannum}
\usepackage{amsfonts}
\usepackage{amssymb}
\usepackage{fancyhdr}
\usepackage{graphicx}
\usepackage{t1enc}
\usepackage{svg}
\usepackage[magyar]{babel}
\usepackage[left=2cm,right=2cm,top=2cm,bottom=2cm]{geometry}

\title{K+F Projekt dokumentáció}
\author{}
\date{2}
\pagestyle{fancy}
\lhead{K+F Projekt dokumentáció}
\rhead{ACSG Kft.}
\cfoot{\thepage. oldal}

\begin{document}
\pagenumbering{arabic}

\begin{center}
    \hrule
    \vspace{0.5cm}
    \begin{Huge}
    K+F projekt dokumentáció\\
    ACSG Kft.\\
    \end{Huge}
    \vspace{0.5cm}
    \begin{huge}
    Időszak:\\
    \end{huge}
    \begin{Large}
    2023.01.01. - 2023.02.28.\\
    \end{Large} \vspace{10pt}
    \begin{large}
    Készítette: Wenesz Dominik\\
    \end{large}\vspace{0.5cm}
    \hrule
    
  \end{center}
  \begin{figure}[b]
    \centering
    \includegraphics[]{acsg.png}
\end{figure}
  


\thispagestyle{empty}
\setcounter{page}{0}




\newpage
\tableofcontents
\newpage

\section{Összefoglaló}
A dokumentum által tekintett időintervallumon belül megtörtént a robotkar (SCARA)
pontosságának tesztelése a munkatérben, melynek eredményeként elmondható, hogy 
a csatlakozóházak megfelelően precíz manipulálásához megfelelő az alkalmazott robotkar.\\[5pt]
A pontosság mérésének tükrében a későbbi digitális iker tervezése is megkezdődött, de 
a tekintett időszakban mindösszesen a szimulációs keretrendszer beállítása indult el 
egy Linux operációs rendszerű PC-n.\\[5pt]
A mélytanulásos objektumdetektálás tanítóadat generáló szoftvere a jelenlegi teszt 
alkalmazásokhoz teljesen elkészült, a későbbiekben a rendszerszintű integráció 
kapcsán lesz csak szükséges további változtatás eszközölése rajta.\\[5pt]
Folytattuk a különböző Deep Learningen alapuló objektum detektáló hálóarchitektúra 
vizsgálatát a csatlakozóházak detektálására.\\[5pt]
Az objektum detektáláson túl a pozíciómeghatározás lehetséges megoldásait vizsgáltuk,
melynek eredményei lentebb láthatók, illetve az objektum detektálást részletesen tárgyaló
dokumentumba felvezetésre kerültek.


\section{Objektum detektáló alrendszer}
Az előző dokumentumban kifejtésre került, hogy a deep learning egy jó alternatívát jelent a hagyományos képfeldolgozási 
módszerekre többek közt az objektumdetektálás területén belül is. 

\subsection{Vizsgált hálóarchitektúrák}

\section{Tanítóadat generálás}
\subsection{Probléma ismertetése}

\subsection{Megoldási lehetőségek}

\subsection{Vegyes tanító adat}

\subsection{Valós adatra illesztés}

\subsection{További fejlesztési irányok}





\section{Pozíciómeghatározás}
A fentebb említett hálóarchitektúrák önmagukban mindösszesen az objektumdetektálásra képesek, mint azt az előző
részekben megállapítottuk, jelen alkalmazási formára a sebesség igényénél fogva elegendő a bounding-box típusú
objektu felismerést alkalmaznunk, hiszen elegendő adatot tartalmaz a robotvezérléshes és szignifikánsan gyorsabb,
mint a szegmentálással is kiegészített hálóarchitektúra.\\[5pt]
Azonban a megfelelő fészekbe való pozícionáláshoz további input paraméterekre van szükségünk, az objektum elfordultsága
is nélkülözhetetlen paraméter ha egy kamerával szeretnénk megoldani a feladatot. Azonban itt több lehetőségünk is van a 
megvalósításra, a dokumentum által tárgyalt időszakban ezen lehetőségeket térképeztük fel és alkalmazhatóságukat 
az adott problémán elkezdtük megvizsgálni.\\
Az általunk megfontolandónak talált megközelítések a következők:
\begin{itemize}
  \item Egy a munkaterület felett elhelyezett kamerától különálló kamera, mely a fészek felett van elhelyezve
  és hagyományos képfeldolgozási technikával (pl.: szögkülönbség meghatározás) megfelelő rotációs értéket adja a 
  robotvezérlőnek. Természetesen léteznek nem explicit szögmeghatározó algoritmusok is, melyek egy adott hibaminimalizációra,
  vagy mintakövetésre törekszenek, azaz kvázi regresszív módszerrel közelítik a megfelelő orientációt, mellyel a fészekbe
  kell kerülnie a csatlakozóháznak.\\
  Ez a megoldás viszonylag egyszerű, azonban mivel két kamera szükséges hozzá (munkatér nagysága miatt is), illetve 
  nem minden csatlakozó esetén feltétlenül hibamentesen generalizálható, ezért alkalmazhatóságát meg kell vizsgálni részletesebben.
  \item Az előző esetet egy kamerával is megoldhatjuk ha azt a robotkar és a munkatér, illetve a fészek geometriája
  lehetővé teszi.
  \item A hagyományos képfeldolgozási módszereken túl ebben az esetben is használhatunk neurális hálózatokat, azaz
  deep learingen alapuló módszereket.\\
  Egyik módszer, hogy már az objektumdetektáló neurális hálózatot bővítjük ki egy orientációt meghatározó résszel. 
  Ekkor annyi szabadságunk van, hogy tetszőleges kondíciókat adhatunk meg, azaz lehetséges a fészekbe való behelyezés
  több variációval való kiegészítése, illetőleg regressziós módon megadhatunk 1-3 elfordulási komponenst (tengelyek körül),
  de tekintve, hogy SCARA robottal dolgozunk, így elegendő egy elfordulással operálni ha már megállapítottuk az 
  objektumról, hogy felvehető.\\
  A rendszer ilyen szintű bővítésével eljutunk oda, hogy implicit meghatározzuk az objektum 3D térben vett 6 szabadságfokú
  pozícióját, mely a mélytanulás ilyen típusú alkalmazásában egy teljesen külön álló kutatási terület, melyre 
  léteznek megoldási módok, melyeket a későbbiekben tesztelni is fogunk, azonban jelen alkalmazásnál SCARA robotnál 
  egyszerűbb modellel is képesek vagyunk elérni a mefelelő pozícionálást elméleti szinten.
  \item A pozíció meghatározása nem feltétlen kell, hogy explicit módon történjen, ha egy adott objektum osztályon 
  belül úgynevezett keypoint-okat keresünk, akkor azok elhelyezkedéséből determinálható az objektum pozíciója.
  Ennek előnye, hogy rengeteg információt tudunk meghatározni egy lépésben, de jellemzően egy ilyen háló betanítása
  több tanítóadatot igényel és nagyobb osztályszám esetén nem feltétlenül konvergál a valósághoz, továbbá természetesen
  a végrehajtási idő is jeentősen megnő a háló komplexitásával együtt.
\end{itemize}
\subsubsection{Megvizsgált módszerek}
Az adott időintervallumban egyelőre főként csak a lehetséges megoldások elméleti kidolgozása zajlott, illetve
a hagyományos képfeldolgozáson alapuló módszerek közül elkezdtük az algoritmusok implementációját is, persze
az OpenCV-ben találhatók hasonló algoritmusok, de a feladat specifikussága miatt mindenképp akár ezeket felhasználva is, 
de saját algoritmust fejlesztenünk a robosztusabb megoldáshoz.
\section{Robot manipulátor}
Miután már kiválasztásra került a technológiai szempontból optimális végberendezés, annak valódi környezetben 
való tesztelését végeztük, továbbá magának a robotvezérlőnek a beállítása, kalibrálása, a robot felprogramozása és 
egyszerűbb illetve összetettebb mozgások megvalósítása LUA programnyelven.\\[5pt]
Többek között a maximális csuklógyorsulást és csuklósebességet teszteltük. Ezek a robot adatlapjában is meg vannak adva 
számszerűen, azonban mindig érdemes az ilyen jellegű validálás, főként mivel az adatlapok az ideális, adott modellezési 
módhoz tartozó értékeket tartalmazzák.\\
(...)

\section{Irodalomjegyzék}







\end{document}